\makereporttitle{2024}{august}

\makereportpreamble

\textbf{Cursuri și seminare științifice:}

În luna august 2024, am început să parcurg cursul de \emph{Advanced Statistical Physics} de pe platforma \href{https://edx.org}{edx.org}, ținut de prof.~Tobias Kippenberg de la \emph{École Polytechnique Fédérale de Lausanne}. Cursul introduce formal noțiuni cum ar fi mișcarea Browniană, ecuația Chapman-Kolmogorov, procese stocastice, calcul stocastic Ito și alte concepte utile.

Am continuat să citesc lucrări de cercetare pe zona de haos cuantic, pentru a înțelege mai bine ce fel de întrebări științifice și ce fel de probleme se mai studiază în acest domeniu.

Între 26 și 29 august am fost la cea de-a 16-a ediție a \href{https://16colfrro.sciencesconf.org/?forward-action=index&forward-controller=index&lang=en}{Colocviului Franco-Român} de matematică aplicată. Am participat la mai multe prelegeri, printre care:

\begin{itemize}
    \item \emph{„Fractional diffusion for kinetic equations”} de Majorlane Puel

    \item \emph{„Localization of energies in Navier-Stokes models with reaction terms”} de Radu Precup

    \item \emph{„Spectral Phases of the Erdős-Rényi graph”} de Johannes Alt

    \item \emph{„Some recent results on the discontinuous Galerkin approximation of Maxwell’s equations”} de Alexandre Ern

    \item \emph{„A dissipative nonlinear Schrodinger model for wave propagation in the marginal ice zone”} de Emilian Părâu

    \item \emph{„Stability in shape optimization under convexity constraint”} de Jimmy Lamboley

    \item \emph{„Maps and 2-dimensional random planar geometries”} de Grégory Miermont

    \item \emph{„Numerical scheme for SPDEs with singular drift”} de Ludovic Goudenge
    
    \item \emph{„Elliptic problems with ‘log-normal’ random coefficients”} de Victor Nistor
    
    \item \emph{„Asymptotic behaviour of a one-dimensional avalanche model through a non-conservative coagulationfragmentation equation and associated stochastic processes”} de Oana Lupașcu-Stamate

    \item \emph{„Numerical spectral analysis of Cauchy-type inverse problems: A probabilistic approach”}

    \item \emph{„Uncertainty principles and sharp stability of functional inequalities”} de Cristian Cazacu
\end{itemize}

\vspace{0.5em}

\textbf{Observațiile conducătorului științific}:

\vspace{3em}

\makereportpostamble{29 august 2024}
